\documentclass[11pt]{article}
\usepackage[a4paper,margin=1in]{geometry}
\usepackage{amsmath,amssymb,amsthm,mathtools}
\usepackage{enumitem}

\newtheorem{theorem}{Theorem}[section]
\newtheorem{lemma}[theorem]{Lemma}
\newtheorem{proposition}[theorem]{Proposition}
\newtheorem{remark}[theorem]{Remark}
\newtheorem{example}[theorem]{Example}

\newcommand{\R}{\mathbb R}
\newcommand{\Om}{\Omega}
\newcommand{\dT}{\delta_T}
\newcommand{\DI}{D_I}
\newcommand{\A}{\mathcal A}
\newcommand{\HH}{\mathcal H}

\title{Does Lipschitz Boundary plus Small \(D_I\) Give Graph Entry?}
\author{}
\date{May 2026}

\begin{document}
\maketitle

\section{Question}

The proposed input is:
\[
   \partial E\hbox{ is Lipschitz},\qquad
   |\nabla u_E|\hbox{ is bounded},\qquad
   \DI(E):=P(E)^2-4\pi |E|\ll1 .
\]
Can one conclude that \(E\) is a small radial graph over a circle?

The answer is:

\begin{quote}
\emph{Not literally from these hypotheses alone.  But small \(D_I\) gives a
strong planar reduction: after discarding satellite components of
quadratic area, the convex hull of the main component is a small-amplitude
radial graph.  To prove that the original set itself is a graph, or that it
can be radialized at the sharp torsion cost, one still needs a no-fold/no-fjord
lemma.}
\end{quote}

\section{Relevant Known Results}

The closest standard results have an extra hypothesis beyond Lipschitz
regularity.

\begin{enumerate}[label=\arabic*.]
\item \textbf{Bonnesen in the plane.}  For a convex planar set \(K\),
\[
   P(K)^2-4\pi |K|\ge \pi^2(R_K-r_K)^2.
\]
Thus a convex set with small isoperimetric deficit lies in a thin annulus and
is a small-amplitude radial graph.

\item \textbf{Fuglede.}  For convex bodies, or for already nearly spherical
sets whose boundary is a small Lipschitz normal graph over the sphere, the
isoperimetric deficit controls the graph norm.  This is a stability theorem
\emph{inside} the convex or graph class; it is not an entry theorem from an
arbitrary Lipschitz boundary into the graph class.

\item \textbf{Cicalese--Leonardi selection.}  The selection principle for the
sharp quantitative isoperimetric inequality creates a better competitor with
quasiminimality/regularity structure.  Again, the regularity comes from the
selected variational problem, not from small perimeter deficit alone.
\end{enumerate}

So the available theorem is not ``Lipschitz plus small \(D_I\) implies graph''.
The available theorem is closer to ``small \(D_I\) plus convexity or
near-spherical parametrization gives a quantitative graph estimate''.  In our
setting, convexification is therefore the natural way to use \(D_I\).

\section{Why the literal statement is false}

\begin{example}[Tiny satellite]
Let \(0<r\ll1\), choose \(R_r=(1-r^2)^{1/2}\), and set
\[
   E_r=B_{R_r}(0)\cup B_r(x_0),
\]
where the two disks are disjoint and \(x_0\) stays in a fixed bounded region.
Then \(|E_r|=\pi\) and
\[
   P(E_r)=2\pi R_r+2\pi r
      =2\pi+2\pi r+O(r^2),
\]
so
\[
   \DI(E_r)=P(E_r)^2-4\pi^2=O(r)\to0 .
\]
The boundary is smooth and uniformly Lipschitz in local charts.  The torsion
gradient is bounded uniformly, since on a disk of radius \(s\) the boundary
gradient is \(s/2\).  But \(E_r\) is not a radial graph over one circle.
\end{example}

Thus connectedness or a rule for deleting negligible components is necessary.
This is not a technicality: small \(D_I\) permits small components, and
Lipschitz regularity does not remove them.

\begin{remark}
Even for connected sets, small local Lipschitz constant and small perimeter
deficit do not by themselves forbid a shallow overhang inside a thin annulus.
Such a set may have several intersections with a ray while paying only a small
extra perimeter.  Therefore small \(D_I\) should be expected to give a
graph envelope or a repair theorem, not literal one-crossing for free.
\end{remark}

\section{What small \(D_I\) does prove in the plane}

Normalize \(|E|=\pi\) and write
\[
   P(E)\le 2\pi+\varepsilon.
\]
Here \(\varepsilon\) is equivalent to \(D_I(E)\) in the small regime:
\[
   \DI(E)=P(E)^2-4\pi^2
      =(P(E)-2\pi)(P(E)+2\pi)
      \simeq 4\pi\varepsilon .
\]

\begin{lemma}[Satellite components are quadratic]
\label{lem:sat}
Let \(E\subset\R^2\) have finite perimeter, \(|E|=\pi\), and
\(P(E)\le2\pi+\varepsilon\), with \(\varepsilon\) small.  Let \(E_0\) be an
indecomposable component of largest area.  Then
\[
   |E\setminus E_0|\le C\varepsilon^2 .
\]
\end{lemma}

\begin{proof}
Let \(a_j\) be the areas of the indecomposable components, with \(a_0\)
largest, and set \(v=\sum_{j\ge1}a_j\).  By the planar isoperimetric
inequality and additivity of perimeter over indecomposable components,
\[
   2\pi+\varepsilon
      \ge P(E)
      \ge 2\sqrt\pi\left(\sqrt{\pi-v}+\sum_{j\ge1}\sqrt{a_j}\right).
\]
Since \(2\sqrt\pi\sqrt{\pi-v}\ge 2\pi-Cv\) for small \(v\),
\[
   \sum_{j\ge1}\sqrt{a_j}\le C\varepsilon+Cv.
\]
Using \(v\le(\sum_{j\ge1}\sqrt{a_j})^2\) and absorbing for small
\(\varepsilon\) gives \(v\le C\varepsilon^2\).
\end{proof}

\begin{lemma}[Convex graph envelope]
\label{lem:convex-envelope}
Let \(F\subset\R^2\) be the largest component from
Lemma~\ref{lem:sat}, and assume \(F\) is a continuum.  Let
\[
   K=\operatorname{co}(F).
\]
Then \(P(K)\le P(F)\le2\pi+C\varepsilon\), and after translating and
rescaling \(K\) to area \(\pi\), one has
\[
   K^\sharp=\{z+re^{i\theta}:0<r<1+\varphi_K(\theta)\},
   \qquad
   \|\varphi_K\|_{L^\infty}\le C\sqrt{\varepsilon}.
\]
\end{lemma}

\begin{proof}
For planar continua, perimeter does not increase under convexification:
\[
   P(\operatorname{co}(F))\le P(F).
\]
The area of \(K\) is bounded and close to \(\pi\) up to a harmless
\(O(\varepsilon)\) error in the small-deficit regime.  Bonnesen's inequality
for convex planar sets gives
\[
   P(K)^2-4\pi |K|
      \ge \pi^2\,(R_K-r_K)^2,
\]
where \(R_K\) and \(r_K\) are the circumradius and inradius of \(K\) with
respect to a suitable center.  Since the left side is \(O(\varepsilon)\),
\[
   R_K-r_K\le C\sqrt\varepsilon .
\]
After area normalization, this is exactly containment in an annulus of
thickness \(C\sqrt\varepsilon\).  A convex body containing the center as an
interior point is a radial graph, and the annular bound gives the stated
\(L^\infty\) amplitude estimate.
\end{proof}

\section{Where Lipschitz and \(|\nabla u|\) help}

The previous two lemmas did not use Lipschitz boundary or the gradient bound.
Those hypotheses are still useful, but they play a different role.

\begin{enumerate}[label=\arabic*.]
\item A uniform Lipschitz boundary gives local finite-overlap charts and
prevents arbitrarily wild fractal boundary behavior.
\item A uniform bound for \(|\nabla u_E|\) prevents high-gradient boundary
defects, such as tiny holes for the torsion problem.
\item In the selected-level setting, the stronger estimate
\[
   \bigl\||\nabla u_E|-c\bigr\|_{L^\infty(\partial E)}\ll1
\]
is much more powerful than mere boundedness.  It rules out small components
and low-gradient tentacles that small \(D_I\) alone permits.
\end{enumerate}

But these hypotheses still do not give the missing radial projection
injectivity.  Local Lipschitz graphs can fold relative to the radial
projection.

\section{Correct theorem to prove}

The useful replacement for literal graph entry is the following.

\begin{proposition}[What the current inputs reduce to]
Assume the stopped planar level \(E\) satisfies
\[
   |E|=\pi,\qquad
   P(E)\le2\pi+\varepsilon,\qquad
   \partial E\hbox{ uniformly Lipschitz},
\]
and has the selected-level torsion bounds
\[
   \|\nabla u_E\|_{L^\infty}\le L,\qquad
   \bigl\||\nabla u_E|-c\bigr\|_{L^\infty(\partial E)}\le\eta.
\]
Then small \(D_I\) gives, quantitatively, a small-amplitude convex radial
graph envelope \(K^\sharp\) for the main component of \(E\), after deleting
satellites of area \(O(\varepsilon^2)\).
\end{proposition}

What remains to close the hybrid route is the sharp radialization estimate
\[
   |E\Delta G|\le C\sqrt{\delta},
   \qquad
   \dT(G)\le C\delta,
\]
where \(G\) is a small-amplitude radial graph.  The new work is to prove that
the parts of \(K\setminus E\) and \(E\setminus K\) coming from folds, fjords,
and outward fingers can be filled or pruned at this torsion cost.  This is
where the \(D_H\) thin-finger estimate and the capacity/no-fjord mechanism
must enter.

\section{Verdict}

There is a standard route from small \(D_I\) to a graph envelope:
component cleanup, convex hull, and Bonnesen.  It is fully quantitative in
dimension two and gives an \(L^\infty\) radial graph after convexification.

There is not, however, a theorem saying that a uniformly Lipschitz set with
bounded torsion gradient and small \(D_I\) is itself a radial graph.  The
literal statement is false without at least deleting negligible components,
and for connected sets it still needs a no-fold/no-fjord argument.  Thus the
right next lemma is not ``\(D_I\) small implies graph''; it is:
\[
   \hbox{convex graph envelope}
   +\hbox{Serrin/}D_H\hbox{ control}
   \Longrightarrow
   \hbox{sharp radialization}.
\]

\end{document}
